% !PRINTED 135
\weberpage{135}
%<<CONT
Theiler von $P$ vom Grade $p^{\beta} q$, wobei $0 \leqq \beta < \alpha$ vorausgesetzt ist, und es sei $\beta$ so gross als möglich angenommen. Dann ist wegen der vorausgesetzten Einfachheit der Gruppe $P$ der Satz \S.\,28, 2. anwendbar, wonach $P$ isomorph ist mit einer Permutationsgruppe von $p^{\alpha-\beta}$ Ziffern. Unter den Permutationen dieser Gruppe giebt es auch solche, deren Grad durch $q$ theilbar ist, und die also, in ihre Cyklen zerlegt, einen Cyklus von $q$ Ziffern enthalten müssen, und daraus folgt:

\weq{1}{p^{\alpha-\beta} > q.}

$\beta)$ Es sei nun zweitens $Q$ ein Theiler von $P$ vom Grade $p^{\alpha}$ und Index $q$. Da nach der Voraussetzung $P$ einfach ist, so kann die in \S.\,30 mit $R$ bezeichnete Gruppe, die aus den mit $Q$ vertauschbaren Elementen von $P$ besteht, nicht mit $P$ identisch sein, da sonst $Q$ ein Normaltheiler von $P$ wäre, was doch, da $P$ als einfache Gruppe vorausgesetzt ist, unmöglich ist. Da aber $Q$ ein Theiler von $R$ ist, und der Index von $Q$ in Bezug auf $P$ eine Primzahl, so muss $R = Q$ sein. Nach II, \S.\,30 giebt es also $q$ und nicht mehr von einander verschiedene conjugirte Gruppen

\[ Q,\ Q_1,\ Q_2,\ \dots\ Q_{q-1}, \]

die alle vom Grade $p^{\alpha}$ sind. Unter diesen Gruppen nehmen wir zwei, etwa $Q$, $Q_1$, die einen grössten gemeinschaftlichen Theiler $Q_0$ vom Grade $p^{\gamma}$ haben, und wir nehmen diese beiden Gruppen so gewählt an, dass $\gamma$ \gesperrt{so gross als möglich} wird. Es ist dann $0 \leqq \gamma < \alpha$. Wenden wir nun auf die Gruppe $P$ das Theorem \S.\,28, 3. an, indem wir für die beiden Gruppen $Q$, $R$ jenes Theorems die Gruppe $Q$ setzen, so ist

\[ q = h_1 + h_2 + h_3 + \cdots, \]

und darin ist $h_1 = 1$, und die übrigen Summanden $h_2, h_3, \dots$ sind die Indices von grössten gemeinschaftlichen Theilern von $Q$ mit je einer der Gruppen $Q_1, Q_2, \dots$. Also sind $h_1, h_2, \dots$ durch $p^{\alpha-\gamma}$ theilbar, % [sic: Weber includes h_1 among the divisible terms although h_1 = 1.]
und es folgt

\[ q \equiv 1 \pmod{p^{\alpha-\gamma}}, \]

also

\weq{2}{q > p^{\alpha-\gamma},}

und folglich wegen (1):

\weq{3}{\gamma > \beta,\quad \gamma > 0.}

$\gamma)$ Ist nun also $Q_0$ der Durchschnitt von $Q$ und $Q_1$ vom Grade $\gamma$, % [sic: Weber prints gamma although Q_0 was defined to have degree p^gamma.]
so können wir nach \S.\,31, III. einen Theiler $R$ von $Q$ vom Grade $p^{\gamma+1}$ bestimmen, von dem $Q_0$ ein Normaltheiler ist, und $R$
%CONT>>
